% Part: first-order-logic
% Chapter: beyond
% Section: intuitionistic-logic

\documentclass[../../../include/open-logic-section]{subfiles}

\begin{document}

\olfileid{fol}{byd}{il}

\olsection{المنطق الحدسي}

على خلاف منطق الرتبة الثانية ومنطق الرتب العليا، يمثل منطق الرتبة
الأولى الحدسي تقييدًا للصورة الكلاسيكية، ويقصد به نمذجة نوع أكثر
«بنائية» من الاستدلال. ولعل الأمثلة الآتية توضح بعض الدوافع الكامنة
وراءه.

افترض أن شخصًا أتاك يومًا وأعلن أنه عيّن عددًا طبيعيًا~$x$ له الخاصية
الآتية: إذا كان $x$ أوليًا كانت فرضية ريمان صادقة، وإذا كان $x$ مركبًا
كانت فرضية ريمان كاذبة. يا لها من بشرى!{} فصدق فرضية ريمان أو كذبها
واحد من أهم الأسئلة المفتوحة في الرياضيات، وها هو ذا قد رد المسألة،
فيما يبدو، إلى عملية حسابية، أي إلى تعيين ما إذا كان عدد بعينه أوليًا
أم لا.

فما قيمة~$x$ السحرية؟ يصفها كما يأتي: $x$ هو العدد الطبيعي المساوي
لـ~$7$ إذا كانت فرضية ريمان صادقة، ولـ~$9$ فيما عدا ذلك.

فتغضب وتطالبه برد مالك. والوصف أعلاه يعيّن فعلًا، من وجهة النظر
الكلاسيكية، قيمة وحيدة لـ~$x$؛ لكن ما تريده حقًا هو قيمة لـ~$x$ معطاة
\emph{صراحة}.

ولنأخذ مثالًا آخر لعله أقل تكلفًا، وانظر في السؤال الآتي. نعلم أنه
يمكن رفع عدد غير نسبي إلى قوة نسبية والحصول على نتيجة نسبية؛ فمثلًا
$\sqrt{2}^2 = 2$. أما إمكان رفع عدد غير نسبي إلى قوة \emph{غير
نسبية} والحصول على نتيجة نسبية فأقل وضوحًا. وتجيب المبرهنة الآتية عن
هذا السؤال بالإيجاب:

\begin{thm}
توجد أعداد غير نسبية $a$ و$b$ يكون $a^b$ نسبيًا.
\end{thm}

\begin{proof}
انظر في $\sqrt{2}^{\sqrt{2}}$. إذا كان هذا العدد نسبيًا، فقد انتهينا:
يمكننا أن نأخذ $a = b = \sqrt{2}$. وإلا فهو غير نسبي. وعندئذ
\[
(\sqrt{2}^{\sqrt{2}})^{\sqrt{2}} = \sqrt{2}^{\sqrt{2} \cdot
  \sqrt{2}} = \sqrt{2}^2 = 2,
\]
وهو نسبي بلا شك. ففي هذه الحالة نأخذ $a$ مساويًا
لـ~$\sqrt{2}^{\sqrt{2}}$، و$b$ مساويًا لـ~$\sqrt 2$.
\end{proof}

هل هذا برهان صحيح؟ يرى معظم الرياضيين أنه كذلك. لكن فيه مرة أخرى ما
لا يبعث على الرضا تمامًا: فقد أثبتنا وجود زوج من الأعداد الحقيقية له
خاصية معينة، من غير أن نستطيع قول \emph{أي} زوج هو. ويمكن إثبات
النتيجة نفسها على نحو يُعطى فيه الزوج $a$، $b$ \emph{داخل} البرهان:
خذ $a = \sqrt{3}$ و$b = \log_3 4$. عندئذ
\[
a^b = \sqrt{3}^{\log_3 4} = 3^{1/2 \cdot \log_3 4} = (3^{\log_3
  4})^{1/2} = 4^{1/2}= 2,
\]
لأن $3^{\log_3 x} = x$. كما أن $b$ غير نسبي: فلو كان
$\log_3 4=p/q$ لعددين صحيحين موجبين، لكان $3^p=4^q=2^{2q}$، وهذا
مستحيل بتفرد التحليل إلى عوامل أولية.

صُمم المنطق الحدسي لنمذجة نوع من الاستدلال لا تجاز فيه خطوات كالخطوة
الواردة في البرهان الأول. فإثبات وجود~$x$ يحقق~$!A(x)$ يعني أن عليك
إعطاء~$x$ بعينه وبرهان أنه يحقق~$!A$، كما في البرهان الثاني. وإثبات
أن $!A$ أو $!B$ صادقة يقتضي أن تستطيع إثبات إحداهما.

وبالعبارة الصورية، منطق الرتبة الأولى الحدسي هو ما نحصل عليه عندما
نقيد نسق !!{derivation} لمنطق الرتبة الأولى على نحو معين. وثمة بالمثل
صور حدسية من منطق الرتبة الثانية ومنطق الرتب العليا. ومن الوجهة
الرياضية ليست هذه إلا أنساقًا استنباطية صورية، لكنها، كما تقدم، تقصد
إلى نمذجة نوع من الاستدلال الرياضي. وقد نعده نوع الاستدلال الذي يسوغه
موقف فلسفي معين من الرياضيات (مثل حدسية براور)، أو نوعًا من الاستدلال
الرياضي أشد «عيانية» وإرضاء (على طريقة بنائية بيشوب)، وقد نختلف في أن
الوصف الصوري يجسد الدافع غير الصوري أو لا. لكن أيًا كانت مواقفنا
الفلسفية، نستطيع دراسة المنطق الحدسي بوصفه منطقًا مقدمًا صوريًا؛
ولأسباب شتى يجد كثير من المناطقة الرياضيين في دراسته فائدة.

للروابط الحدسية تأويل بنائي غير صوري يُعرف عادة بتأويل BHK (نسبة إلى
براور وهايتنغ وكولموغوروف). ومضمونه ما يأتي: يتألف برهان $!A \land
!B$ من برهان لـ~$!A$ مقترن ببرهان لـ~$!B$؛ ويتألف برهان $!A \lor !B$
من برهان لـ~$!A$ أو برهان لـ~$!B$، مع معلومة صريحة تعيّن أيهما حاصل؛
ويتألف برهان $!A \lif !B$ من إجراء يحوّل برهانًا لـ~$!A$ إلى برهان
لـ~$!B$؛ ويتألف برهان $\lforall[x][!A(x)]$ من إجراء يعيد برهانًا
لـ~$!A(x)$ لكل قيمة لـ~$x$؛ ويتألف برهان $\lexists[x][!A(x)]$ من قيمة
لـ~$x$ مع برهان أن هذه القيمة تحقق~$!A$. ويمكن وصف هذا التأويل بعبارات
حسابية فيما يعرف بـ«تماثل كاري--هوارد» أو «نموذج ال!!{formula}s
بوصفها أنماطًا»: فكر في !!a{formula} بوصفها تحدد نوعًا معينًا من
أنماط البيانات، وفي البراهين بوصفها كائنات حسابية من تلك الأنماط تبين
لنا صدق ال!!{formula} المقابلة.

يُنظر كثيرًا إلى المنطق الحدسي بوصفه المنطق الكلاسيكي «ناقص» قانون
الثالث المرفوع. وتجعل المبرهنة الآتية هذا القول أدق.
\begin{thm}
مخططات البديهيات الآتية متكافئة حدسيًا:
\begin{enumerate}
\item $(\lnot !A \lif \lfalse) \lif !A$.
\item $!A \lor \lnot !A$
\item $\lnot \lnot !A \lif !A$
\end{enumerate}
\end{thm}
واشتقاق أمثلة أحد المخططات من أي من المخططين الآخرين تمرين حسن في
المنطق الحدسي.

تُنسب أول الأنساق الاستنباطية للمنطق القضوي الحدسي، التي قدمت صيغًا
صورية لحدسية براور، إلى كولموغوروف وغليفينكو وهايتنغ، كل على استقلال.
وتُنسب أول صياغة صورية لمنطق الرتبة الأولى الحدسي (ولأجزاء من
الرياضيات الحدسية) إلى هايتنغ. ومع أن مخططات عديدة صحيحة كلاسيكيًا لا
تصح حدسيًا، فإن مخططات كثيرة أخرى تصح.

تصف \emph{ترجمة النفي المزدوج} علاقة مهمة بين المنطق الكلاسيكي
والمنطق الحدسي. وهي تعرّف استقرائيًا كما يأتي (فكر في~$!A^N$ بوصفها
الترجمة «الحدسية» لل!!{formula} الكلاسيكية~$!A$):
\begin{align*}
!A^N & \ident \lnot\lnot !A \quad \text{لل!!{formula}s الذرية $!A$} \\
(!A \land !B)^N & \ident (!A^N \land !B^N) \\
(!A \lor !B)^N & \ident  \lnot\lnot (!A^N \lor !B^N) \\
(!A \lif !B)^N & \ident (!A^N \lif !B^N) \\
(\lforall[x][!A])^N & \ident \lforall[x][!A^N] \\
(\lexists[x][!A])^N & \ident \lnot\lnot\lexists[x][!A^N]
\end{align*}
وكان لكولموغوروف وغليفينكو صورتان من هذه الترجمة للمنطق القضوي؛ أما
في منطق المحمولات فتُنسب إلى غودل وغنتسن، كل على استقلال. ولدينا

\begin{thm}
\begin{enumerate}
\item $!A \liff !A^N$ قابلة للإثبات كلاسيكيًا.
\item إذا كانت $!A$ قابلة للإثبات كلاسيكيًا، كانت $!A^N$ قابلة
  للإثبات حدسيًا.
\end{enumerate}
\end{thm}

ويمكننا الآن تصور الحوار الآتي. الرياضي الكلاسيكي: «أثبتُ~$!A$!»
الرياضي الحدسي: «برهانك غير صحيح. ما أثبتَّه حقًا هو~$!A^N$.» الرياضي
الكلاسيكي: «لا بأس عندي!» فمن وجهة نظر الرياضي الكلاسيكي لا يفعل
الحدسي سوى التدقيق في فروق لا طائل منها، لأن الصيغتين متكافئتان. لكن
الحدسي يصر على وجود فرق.

لاحظ أن الترجمة أعلاه لا تتعلق إلا بالمنطق الخالص؛ فلا تتناول سؤال
ما البديهيات \emph{غير المنطقية} المناسبة للرياضيات الكلاسيكية
والحدسية، ولا العلاقة بينها. لكن التوسيع الطفيف الآتي للمبرهنة أعلاه
يوفر معلومات مفيدة:

\begin{thm}
إذا كانت $\Gamma$ تثبت $!A$ كلاسيكيًا، كانت $\Gamma^N$ تثبت $!A^N$
حدسيًا.
\end{thm}

أي إنه إذا كانت $!A$ قابلة للإثبات كلاسيكيًا من بعض الفرضيات، كانت
$!A^N$ قابلة للإثبات من ترجمات النفي المزدوج لتلك الفرضيات.

لبيان أن جملة أو !!{formula} قضوية صحيحة حدسيًا لا يلزمك إلا تقديم
برهان. لكن كيف تبين أنها ليست صحيحة؟ نحتاج لهذا الغرض إلى دلالة سليمة،
ويفضل أن تكون تامة. وتفي دلالة وضعها كريبكه بهذا الغرض وفاء حسنًا.

ويمكننا أن نفعل ما فعلناه في المنطق الكلاسيكي: نعرّف الدلالة ونثبت
السلامة والتمام. غير أن من الجدير ملاحظة الفرق الآتي. كانت الدلالة في
المنطق الكلاسيكي هي الدلالة «الواضحة» بمعنى ما، إذ كانت ضمنية في معنى
الروابط. ومع إمكان تقديم بعض الدوافع الحدسية لدلالة كريبكه، فإنها لا
تبعث الشعور نفسه بالحتمية. يضاف إلى ذلك أن مفهوم ال!!{structure}
الكلاسيكية مفهوم رياضي طبيعي؛ فنستطيع إما اتخاذ مفهوم ال!!{structure}
أداة لدراسة منطق الرتبة الأولى الكلاسيكي، وإما اتخاذ منطق الرتبة الأولى
الكلاسيكي أداة لدراسة ال!!{structure}s الرياضية. وعلى النقيض من ذلك لا
يمكن النظر إلى !!{structure}s كريبكه إلا بوصفها بناء منطقيًا؛ فلا يبدو
أن لها أهمية رياضية مستقلة.

تتألف !!{structure} كريبكه~$\mModel M = \tuple{W, R, V}$ للغة قضوية
من مجموعة~$W$، وترتيب جزئي~$R$ على~$W$ له أصغر !!{element}، وإسناد
«رتيب» للمتغيرات القضوية إلى !!{element}s~$W$. والحدس هو أن
!!{element}s~$W$ تمثل «عوالم» أو «حالات معرفة»؛ وأن العنصر $v \geq u$
يمثل «حالة مستقبلية ممكنة» لـ~$u$؛ وأن المتغيرات القضوية المسندة
إلى~$u$ هي القضايا المعروف صدقها في الحالة~$u$. ثم توسع علاقة الإلزام
$\mSat{M}{!A}[w]$ هذه العلاقة إلى ال!!{formula}s الاعتباطية في اللغة؛
واقرأ $\mSat{M}{!A}[w]$ على أنها «$!A$ صادقة في الحالة~$w$». وتعرّف
العلاقة استقرائيًا كما يأتي:
\begin{enumerate}
\item $\mSat{M}{\Obj p_i}[w]$ إذا وفقط إذا كان $\Obj p_i$ واحدًا من
  المتغيرات القضوية المسندة إلى~$w$.
\item $\mSat/{M}{\lfalse}[w]$.
\item $\mSat{M}{(!A \land !B)}[w]$ إذا وفقط إذا كان
  $\mSat{M}{!A}[w]$ و$\mSat{M}{!B}[w]$.
\item $\mSat{M}{(!A \lor !B)}[w]$ إذا وفقط إذا كان
  $\mSat{M}{!A}[w]$ أو $\mSat{M}{!B}[w]$.
\item $\mSat{M}{(!A \lif !B)}[w]$ إذا وفقط إذا كان، متى كان
  $w' \geq w$ و$\mSat{M}{!A}[w']$، أن $\mSat{M}{!B}[w']$.
\end{enumerate}
ومن التمارين الحسنة أن تحاول بيان أن $\lnot (p \land q) \lif
(\lnot p \lor \lnot q)$ ليست صحيحة حدسيًا، بإنشاء !!{structure}
كريبكه تقدم مثالًا مضادًا.

\end{document}
